% !PRINTED 495
\weberpage{495}

From the simultaneous consideration of the decomposition of permutations into cycles and into transpositions there results a proof of the division of the permutations into two classes, with which we have already become acquainted in the second section, in connection with determinants.

If $\tau$ denotes a transposition and $\pi$ an arbitrary permutation, then in the compounded permutation $\tau\pi$ the number of cycles is greater by one or smaller by one than in $\pi$. The two figures that are interchanged by $\tau$ can occur either in the same cycle of $\pi$ or in two different cycles. Let us assume that a cycle occurring in $\pi$ is $\gamma = (1, 2 \dots a, a + 1 \dots b)$, and that $\tau$ contains two figures that occur in $\gamma$, say $\tau = (1, a)$; then
\[
\tau\gamma = (1, a + 1 \dots b)\ (2, 3 \dots a),
\]
that is, $\gamma$ is decomposed by composition with $\tau$ into two cycles, while the remaining cycles are not touched by $\tau$.

But if the two figures of $\tau$ occur in two different cycles of $\pi$, then let these two cycles be
\[
\gamma = (1, 2, 3 \dots a), \quad \gamma' = (1', 2', 3' \dots a')
\]
and $\tau = (1, 1')$. Then
\[
\tau\gamma\gamma' = (1, 2', 3' \dots a', 1', 2, 3 \dots a),
\]
that is, the two cycles $\gamma$, $\gamma'$ are united by $\tau$ into a single cycle, while again the remaining cycles stay unchanged.

If, then, we have represented a permutation $\pi$ of $m$ figures, consisting of $\nu$ cycles (the one-term cycles being counted along with the rest), by $\mu$ transpositions, then by composition with these $\mu$ transpositions in reversed order it can be converted into the identical permutation, which consists of $m$ one-term cycles. In all, therefore, $m - \nu$ cycles have been gained, and since to every transposition there corresponds a gained or a lost cycle, we must have $\mu \equiv m - \nu\ (\mathrm{mod}.\ 2)$; that is, $\mu$ is an even or an odd number according as $m - \nu$ is even or odd. But the last number depends only on the permutation $\pi$, not on the manner of the decomposition into transpositions. We have thereby the theorem:
%CONT>>
